\documentclass[aspectratio=169,12pt]{beamer}
\usetheme{Madrid}
\usecolortheme{seahorse}
\usepackage[utf8]{inputenc}
\usepackage[T1]{fontenc}
\usepackage{amsmath,amssymb}
\usepackage{booktabs}
\usepackage{array}

\newcommand{\MeV}{\,\text{MeV}}
\newcommand{\GeV}{\,\text{GeV}}
\newcommand{\fpi}{f_\pi}
\newcommand{\zch}{\mathfrak{z}}

\setbeamertemplate{navigation symbols}{}
\setbeamertemplate{footline}[frame number]

\title{Three families from SUSY confinement}
\subtitle{Spectrum and mass predictions of an
$\mathcal{N}=1$ $\mathrm{SU}(3)\times\mathrm{SU}(2)$ theory}
\author{A.~Rivero}
\institute{Universidad de Zaragoza}
\date{2026}

\begin{document}

%-------------------------------------------------------------
\begin{frame}
\titlepage
\end{frame}

%-------------------------------------------------------------
\begin{frame}{The flavour problem}

The Standard Model has \textbf{13 free parameters} in the
charged fermion sector:
\begin{itemize}
\item 9 masses: $m_e, m_\mu, m_\tau, m_u, m_d, m_s, m_c, m_b, m_t$
\item 4 CKM angles: $\theta_{12}, \theta_{23}, \theta_{13}, \delta_{\rm CP}$
\end{itemize}

\bigskip
No established principle relates them.

\bigskip
\textbf{This talk:} an $\mathcal{N}=1$ SUSY theory that
reduces 13 to \textbf{6} free parameters.

\end{frame}

%-------------------------------------------------------------
\begin{frame}{The foundational relation}

\begin{equation*}
\boxed{m_\pi^2 - m_\mu^2 \;\approx\; \fpi^2}
\end{equation*}

\bigskip
\begin{center}
\begin{tabular}{lc}
$m_\pi^2 - m_\mu^2$ & $8\,316\MeV^2$ \\
$\fpi^2$ & $8\,501\MeV^2$ \\
Discrepancy & $2.2\%$
\end{tabular}
\end{center}

\bigskip
If the pion and the muon are \textbf{superpartners}
split by soft SUSY breaking, the splitting scale is~$\fpi$.

\end{frame}

%-------------------------------------------------------------
\begin{frame}{Diophantine bootstrap}

\textbf{Postulate:} every scalar of the SSM is a composite
(meson or diquark) of quarks.

\medskip
Matching composites to fermion d.o.f.\ gives three equations:
\begin{align*}
rs &= 2N \\
\tfrac{r(r+1)}{2} &= 2N \\
r^2 + s^2 - 1 &= 4N
\end{align*}

\pause
\textbf{Unique solution:}
\begin{equation*}
\boxed{(N,\,r,\,s) = (3,\,3,\,2)}
\end{equation*}

Three colours, three generations, five light quarks
$\to$ $\mathrm{SU}(5)$ flavour symmetry.

\medskip
Right-handed neutrinos are \emph{required}. Top quark is
automatically excluded.

\end{frame}

%-------------------------------------------------------------
\begin{frame}{$\mathrm{SU}(5)$ composites}

Composites fill:
\begin{equation*}
\mathbf{24} \oplus \mathbf{15} \oplus \overline{\mathbf{15}}
\end{equation*}

\begin{itemize}
\item $\mathbf{24}$: mesons ($q\bar q$, colour singlet)
--- leptons + sleptons
\item $\mathbf{15} + \overline{\mathbf{15}}$:
diquarks ($qq$, colour triplet)
\end{itemize}

\bigskip
\textbf{Anomaly-free} ($A(\mathbf{24})+A(\mathbf{15})+A(\overline{\mathbf{15}})=0$)

\textbf{Asymptotically free} ($b_0 = 31/3$)

\bigskip
The symmetric $\mathbf{15}$ (not the antisymmetric
$\overline{\mathbf{10}}$) is forced by the counting.

\end{frame}

%-------------------------------------------------------------
\begin{frame}{The Lagrangian: Seiberg dual}

$\mathrm{SU}(3)$ SQCD with $N_f = N_c = 3$.
Confined d.o.f.: meson matrix $M^i{}_j$, baryons $B, \bar B$,
Lagrange multiplier $X$.

\begin{equation*}
\begin{split}
W &= \mathrm{Tr}(\hat m\, M)
  + X\!\left(\det M^{(3)} - B\bar B - \Lambda^{6}\right)
  + c_3\,\frac{(\det M^{(3)})^3}{\Lambda^{18}}\\
&\quad + y_c H_u M^c{}_c + y_b H_d M^b{}_b
  + y_t H_u Q^t \bar Q_t
  + \lambda_S S H_u\!\cdot\!H_d + \tfrac{\kappa}{3}S^3
\end{split}
\end{equation*}

\medskip
\textbf{Seiberg seesaw:}
$\langle M_j\rangle = C/m_j$
\quad (heavy quarks $\to$ light mesons)

\end{frame}

%-------------------------------------------------------------
\begin{frame}{O'Raifeartaigh mechanism}

$W = f\Phi_0 + m\Phi_1\Phi_2 + g\Phi_0\Phi_1^2$

\medskip
$F_{\Phi_0} = f \neq 0$ \quad$\Longrightarrow$\quad
SUSY \textbf{necessarily broken}

\bigskip
Fermion spectrum at vacuum $gv/m = \sqrt{3}$:
\begin{equation*}
\bigl(0,\;\; (2{-}\sqrt{3})\,m,\;\; (2{+}\sqrt{3})\,m\bigr)
\end{equation*}

\pause
\textbf{Structural mass ratio:}
\begin{equation*}
R_{\mathrm{OR}} = \frac{m_+}{m_-}
= (2+\sqrt{3})^2 = 7 + 4\sqrt{3} \approx 13.93
\end{equation*}

No free parameters.

\end{frame}

%-------------------------------------------------------------
\begin{frame}{The mass-charge identity}

Write each mass as $m_k = \zch_k^2$ with
$\zch_k = z_0 + z_k$, $\sum z_k = 0$.

\bigskip
At $gv/m = \sqrt{3}$, the O'Raifeartaigh vacuum satisfies:
\begin{equation*}
\boxed{\langle z_k^2\rangle = z_0^2}
\qquad\Longleftrightarrow\qquad
\frac{\sum m_k}{(\sum\sqrt{m_k})^2} = \frac{2}{3}
\end{equation*}

\bigskip
This is a \textbf{theorem} about the superpotential ---
not an empirical fit.

\medskip
\textbf{Not an RG attractor:} one-loop RG drives
$\langle z_k^2\rangle/z_0^2 \to 0$.
It is a \emph{boundary condition}.

\end{frame}

%-------------------------------------------------------------
\begin{frame}{The mass chain: five quarks from one input}

\begin{equation*}
\boxed{
\begin{aligned}
m_u + m_d &= 6.83\MeV \quad\text{(input)}\\[4pt]
&\xrightarrow{\;\text{isospin seed}\;}
m_s = 95.1\MeV \quad[+1.8\%]\\[4pt]
&\xrightarrow{\;R_{\mathrm{OR}} = (2{+}\sqrt{3})^2\;}
m_c = 1301\MeV \quad[+2.0\%]\\[4pt]
&\xrightarrow{\;\sqrt{m_b} = 3\sqrt{m_s} + \sqrt{m_c}\;}
m_b = 4233\MeV \quad[+1.3\%]\\[4pt]
&\xrightarrow{\;\langle z_k^2\rangle = z_0^2\;\text{on }(c,b,t)\;}
m_t = 171{,}250\MeV \quad[-0.9\%]\\[4pt]
&\xrightarrow{\;\text{dual identity on }(d,s,b)\;}
m_d = 4.6\MeV \quad[-1.5\%]
\end{aligned}
}
\end{equation*}

Spans a factor of $2.5 \times 10^4$ in mass.

\end{frame}

%-------------------------------------------------------------
\begin{frame}{The bion mass relation}

Monopole-instantons on $\mathbb{R}^3\times S^1$ (\"Unsal):
each carries a $\sqrt{m_k}$ factor.

\medskip
Bion K\"ahler correction:
$\delta K \sim \lambda_2\,|S_{\rm bloom} - 2\,S_{\rm seed}|^2$

\medskip
\textbf{Perfect square} $\to$ minimum at:
\begin{equation*}
\boxed{\sqrt{m_b} = 3\sqrt{m_s} + \sqrt{m_c}}
\end{equation*}

The coefficient $3 = N_c$.

\bigskip
\begin{center}
\begin{tabular}{lcc}
\toprule
Input $m_c$ & $m_b$ predicted & PDG pull \\
\midrule
1301 MeV (chain) & 4233 MeV & $+1.8\sigma$ \\
1275 MeV (PDG) & 4186 MeV & $+0.2\sigma$ \\
\bottomrule
\end{tabular}
\end{center}

\end{frame}

%-------------------------------------------------------------
\begin{frame}{The dual mass-charge identity}

The Seiberg seesaw $\langle M_j\rangle = C/m_j$ maps
$\zch_k = \sqrt{m_k}$ to dual charges
$\xi_k \propto 1/\sqrt{m_k}$.

\begin{equation*}
\frac{\sum 1/m_k}{(\sum 1/\sqrt{m_k})^2} = \frac{2}{3}
\end{equation*}

\medskip
For $(d,s,b)$: deviation from $2/3$ is only $0.22\%$.

\medskip
Predicts $m_d = 4.605\MeV$ \quad (PDG: $4.67 \pm 0.48$, $-0.13\sigma$)

\bigskip
\textbf{Not an independent assumption:} algebraic property
of the seesaw $F$-term equations.

\end{frame}

%-------------------------------------------------------------
\begin{frame}{Lepton sector: SU(2) $s$-confinement}

Separate $\mathrm{SU}(2)$ SQCD with $N_f = 3$ ($s$-confining).

Same O'Raifeartaigh structure, same $gv/m = \sqrt{3}$.

\bigskip
No bion mechanism in SU(2) $\to$ spectrum stays near seed.

\bigskip
$\langle z_k^2\rangle = z_0^2$ on $(m_e, m_\mu, m_\tau)$:

\begin{equation*}
\boxed{m_\tau^{\rm pred} = 1776.97\MeV}
\end{equation*}

\begin{center}
PDG: $1776.86 \pm 0.12\MeV$ \qquad
\textbf{$+0.006\%$, $\mathbf{0.9\sigma}$}
\end{center}

\medskip
Most precise prediction in the programme.

\end{frame}

%-------------------------------------------------------------
\begin{frame}{Supermultiplet identification}

Each meson superfield $M^i{}_j$: fermion = lepton,
scalar = meson.

\bigskip
\begin{center}
\begin{tabular}{lccc}
\toprule
Multiplet & Fermion & Scalar & SUSY-limit mass \\
\midrule
$\Phi_0$ (goldstino) & $e$ & $\pi$ & $\to 0$ \\
$\Phi_1$ (light) & $\mu$ & $K$ & $(2{-}\sqrt{3})\,m$ \\
$\Phi_2$ (heavy) & $\tau$ & $D_s$ & $(2{+}\sqrt{3})\,m$ \\
\bottomrule
\end{tabular}
\end{center}

\bigskip
Soft mass $\tilde m^2 = \fpi^2$ splits each pair:
\begin{itemize}
\item $m_\pi/m_e \approx 264$ (large splitting, lightest)
\item $m_K/m_\mu \approx 4.7$ (moderate)
\item $m_{D_s}/m_\tau \approx 1.11$ (near-degenerate, heaviest)
\end{itemize}

\end{frame}

%-------------------------------------------------------------
\begin{frame}{Electroweak predictions}

\textbf{$\tan\beta = 1$} is required: mass-charge identity on
masses $=$ identity on Yukawas only when $\sin\beta = \cos\beta$.

\bigskip
\textbf{Higgs mass} (NMSSM, $S \neq X$):
\begin{equation*}
m_h = \frac{\lambda_S v}{\sqrt{2}} = 125.35\GeV
\qquad (\lambda_S = 0.72,\;\; 0.6\sigma)
\end{equation*}

\bigskip
\textbf{Cabibbo angle} (GST from seesaw texture):
\begin{equation*}
\sin\theta_C = \sqrt{m_d/m_s} = 0.2236
\qquad (\text{PDG:}\; 0.2243,\;\; 0.9\sigma)
\end{equation*}

\bigskip
\textbf{Weak mixing} (Casimir eigenvalue equation):
\begin{equation*}
\sin^2\theta_W = 0.22310
\qquad (\text{PDG:}\; 0.22320,\;\; 0.4\sigma)
\end{equation*}

\end{frame}

%-------------------------------------------------------------
\begin{frame}{Summary of predictions}

\small
\begin{center}
\begin{tabular}{llll}
\toprule
Quantity & Predicted & PDG 2024 & Source \\
\midrule
$m_s$ & $95.1\MeV$ & $93.4 \pm 8.6$ ($0.2\sigma$)
  & Isospin seed \\
$m_c$ & $1301\MeV$ & $1275 \pm 25$ ($1.0\sigma$)
  & O'Raifeartaigh \\
$m_b$ & $4233\MeV$ & $4180 \pm 30$ ($1.8\sigma$)
  & Bion K\"ahler \\
$m_t$ & $171.3\GeV$ & $172.76 \pm 0.30$ ($-0.9\%$)
  & Yukawa constraint \\
$m_d$ & $4.6\MeV$ & $4.67 \pm 0.48$ ($0.1\sigma$)
  & Dual identity \\
$m_\tau$ & $1776.97\MeV$ & $1776.86 \pm 0.12$ ($0.9\sigma$)
  & Lepton identity \\
$V_{us}$ & $0.2236$ & $0.2243 \pm 0.0008$ ($0.9\sigma$)
  & GST \\
$m_h$ & $125.35\GeV$ & $125.25 \pm 0.17$ ($0.6\sigma$)
  & NMSSM \\
$\sin^2\theta_W$ & $0.22310$ & $0.22320$ ($0.4\sigma$)
  & Casimir eq. \\
\bottomrule
\end{tabular}
\end{center}

\textbf{All pulls below $2\sigma$.}

\end{frame}

%-------------------------------------------------------------
\begin{frame}{Parameter budget}

\begin{columns}[T]
\column{0.48\textwidth}
\textbf{Seven eliminated:}
\begin{enumerate}
\item $m_s$ from $m_u{+}m_d$ (isospin seed)
\item $m_c$ from $m_s$ (O'Raifeartaigh)
\item $m_b$ from $m_s, m_c$ (bion)
\item $m_t$ from $m_c, m_b$ (Yukawa)
\item $m_d$ from $m_s, m_b$ (dual)
\item $m_\tau$ from $m_e, m_\mu$ (lepton)
\item $\theta_{12}$ from $m_d, m_s$ (GST)
\end{enumerate}

\column{0.48\textwidth}
\textbf{Six remain free:}
\begin{itemize}
\item $m_u$
\item $m_e$
\item $m_\mu$
\item $\theta_{23}$
\item $\theta_{13}$
\item $\delta_{\rm CP}$
\end{itemize}

\bigskip
$\chi^2/\text{dof} = 0.12$

$p = 0.97$
\end{columns}

\end{frame}

%-------------------------------------------------------------
\begin{frame}{Spectrum flow}

\small
\begin{center}
\begin{tabular}{p{2.8cm}p{3.5cm}p{5cm}l}
\toprule
Stage & Fermions & Scalars & STr \\
\midrule
Full SUSY & $(0, m_-, m_+)$
  & exact boson--fermion pairs & $0$ \\[3pt]
$F$-breaking & unchanged
  & $B$-term: $\pi < m_k^2 < \sigma$ & $0$ \\[3pt]
CW stabilisation & unchanged
  & pseudo-modulus lifted & $0$ \\[3pt]
Bion $K$ & $v_0$-doubling
  & scalars track & $0$ \\[3pt]
3-instanton & $\delta$ selected
  & angular potential & $0$ \\[3pt]
Soft $\fpi^2$ & unchanged
  & mesons at $\sqrt{2}\fpi$ & $18\fpi^2$ \\
\bottomrule
\end{tabular}
\end{center}

STr $= 0$ at every stage except the last.

\end{frame}

%-------------------------------------------------------------
\begin{frame}{The flavour-universal Yukawa}

Seiberg seesaw: $\langle M_j\rangle = C/m_j$

Standard Yukawa: $y_j = \sqrt{2}\,m_j/v_{\rm EW}$

\bigskip
Their product:
\begin{equation*}
y_j\,\langle M_j\rangle
= \frac{\sqrt{2}\,m_j}{v_{\rm EW}}\,\frac{C}{m_j}
= \frac{\sqrt{2}\,C}{v_{\rm EW}}
= 5.07\MeV
\end{equation*}

\textbf{Same for every confined flavour.} The $m_j$ cancels exactly.

\bigskip
$\Longrightarrow$ Glashow--Weinberg FCNC suppression is
\emph{automatic}.

$\Longrightarrow$ Yukawa couplings are \emph{not free parameters}.

\end{frame}

%-------------------------------------------------------------
\begin{frame}{The $\mathrm{SO}(32)$ embedding}

$\mathrm{SO}(32) \supset \mathrm{SU}(16) \supset
\mathrm{SU}(5)_f \times \mathrm{SU}(3)_c$

\medskip
$\mathbf{16} = (\mathbf{5},\mathbf{3}) \oplus
(\mathbf{1},\mathbf{1})$
\qquad ($5 \times 3 + 1 = 16$)

\bigskip
$\mathbf{496} = \wedge^2(\mathbf{32})$ contains:

\begin{center}
\begin{tabular}{ll}
$(\mathbf{24},\mathbf{1})$ & mesons \\
$(\mathbf{15},\bar{\mathbf{3}})$ & diquarks \\
$(\overline{\mathbf{15}},\mathbf{3})$ & anti-diquarks
\end{tabular}
\end{center}

\bigskip
The symmetric $\mathbf{15}$ appears naturally from
$S^2(\mathbf{5}) \otimes \wedge^2(\mathbf{3})$.

\medskip
$E_8 \times E_8$ is \textbf{blocked}: only exterior powers
of $\mathbf{5}$, never $S^2(\mathbf{5})$.

\end{frame}

%-------------------------------------------------------------
\begin{frame}{Open problems}

\begin{enumerate}
\item \textbf{Diquark obstruction:} Pauli theorem forces
antisymmetric $\overline{\mathbf{10}}$ for $qq$ composites,
but the counting requires symmetric $\mathbf{15}$.
The diquark must be non-composite.

\bigskip
\item \textbf{Bloom mechanism:} The seed $gv/m = \sqrt{3}$ is
an unstable fixed point under $R$-breaking. The instability
points in the right direction, but the bloom magnitude is
nonperturbative.

\bigskip
\item \textbf{$V_{cb}$:} Fritzsch 6-zero texture gives
$|V_{cb}| \geq 0.059$ (PDG: $0.042$, 40\% too large).
The 2--3 sector needs a modified texture or separate mechanism.

\bigskip
\item \textbf{K\"ahler stabilisation:} Both $V_{\rm tree}$ and
$V_{\rm CW}$ are monotonically increasing --- the pole is a wall,
not a well. The minimum at $\sqrt{3}$ requires a nonperturbative
contribution.
\end{enumerate}

\end{frame}

%-------------------------------------------------------------
\begin{frame}{Conclusions}

\begin{itemize}
\item SUSY + confinement uniquely select $N = 3$ generations
  via a Diophantine system

\bigskip
\item The O'Raifeartaigh superpotential at $gv/m = \sqrt{3}$
  produces the mass-charge identity as a \emph{theorem}

\bigskip
\item From $m_u + m_d$ alone: five quark masses predicted to
  $\leq 2\%$ accuracy

\bigskip
\item Lepton mass-charge identity: $m_\tau$ to $0.006\%$

\bigskip
\item Seven of thirteen SM parameters eliminated
  ($\chi^2/\text{dof} = 0.12$, $p = 0.97$)

\bigskip
\item The theory lives naturally in $\mathrm{SO}(32)$
  (Type~I string), not $E_8 \times E_8$
\end{itemize}

\end{frame}

%-------------------------------------------------------------
\begin{frame}{}
\centering
\Huge Thank you

\bigskip
\large
\texttt{arXiv:2407.05397}

\bigskip
\normalsize
\textit{Eur.\ Phys.\ J.\ C} \textbf{84}, 1058 (2024)
\end{frame}

%=============================================================
% BACKUP SLIDES
%=============================================================

\appendix
\newcounter{finalframe}
\setcounter{finalframe}{\value{framenumber}}

%-------------------------------------------------------------
\begin{frame}{Backup: the mass-charge identity and Koide}

The identity
$\displaystyle\frac{\sum m_k}{(\sum\sqrt{m_k})^2} = \frac{2}{3}$
was discovered empirically by Y.~Koide (1982) for the
charged leptons $(e,\mu,\tau)$.

\bigskip
\textbf{Key difference in this work:}
\begin{itemize}
\item Koide's formula is an \emph{observation} --- unexplained,
possibly accidental
\item Here it is a \textbf{theorem}: the O'Raifeartaigh
superpotential at $gv/m = \sqrt{3}$ produces
$(0, m_-, m_+)$ satisfying the identity \emph{exactly}
\item The identity is a boundary condition set by the
superpotential, not an RG attractor
\end{itemize}

\bigskip
We avoid the name ``Koide'' in the paper to prevent confusion
with the empirical observation: the identity here has a
Lagrangian derivation.

\end{frame}

%-------------------------------------------------------------
\begin{frame}{Backup: numerical status of the identity}

\small
\begin{center}
\begin{tabular}{lccc}
\toprule
Triple & $\displaystyle\frac{\sum m_k}{(\sum\sqrt{m_k})^2}$
  & Deviation from $2/3$ & Status \\
\midrule
$(e,\mu,\tau)$ & $0.6667$ & $0.006\%$ & exact (input) \\[4pt]
$(0, m_u{+}m_d, m_s)$ & $0.665$ & $0.27\%$ & isospin seed \\[4pt]
$(-\sqrt{m_s},\sqrt{m_c},\sqrt{m_b})$ & $0.674$
  & $1.2\%$ & bloomed \\[4pt]
$(\sqrt{m_c},\sqrt{m_b},\sqrt{m_t})$ & $0.669$
  & $0.4\%$ & Yukawa constraint \\
\bottomrule
\end{tabular}
\end{center}

\bigskip
\begin{itemize}
\item Lepton triple: near-exact (bloom is only $2.3^\circ$)
\item Quark triples: bloomed away from seed, but $< 2\%$
\item The seed ($\delta = 3\pi/4$, one mass zero) satisfies
  the identity \emph{exactly} by construction
\end{itemize}

\end{frame}

%-------------------------------------------------------------
\begin{frame}{Backup: the dual identity}

The Seiberg seesaw $\langle M_j\rangle = C/m_j$ maps
$\sqrt{m_k} \to 1/\sqrt{m_k}$:

\begin{equation*}
\frac{\sum 1/m_k}{(\sum 1/\sqrt{m_k})^2} = \frac{2}{3}
\end{equation*}

\medskip
\begin{center}
\begin{tabular}{lcc}
\toprule
Triple & Direct ratio & Dual ratio \\
\midrule
$(d,s,b)$ & $0.732$ (far) & $0.665$ ($0.22\%$) \\
$(u,d,s)$ & seed (exact) & seed (exact) \\
\bottomrule
\end{tabular}
\end{center}

\bigskip
At the seed (one mass $= 0$), direct and dual coincide.

After bloom, the seesaw \emph{preserves} the identity much
better than the direct masses: it maps the bloomed spectrum
back toward the seed hierarchy.

\medskip
Dual self-consistency (both $= 2/3$): requires
$\cos 3\delta = 5\sqrt{2}/8$.

\end{frame}

%-------------------------------------------------------------
\begin{frame}{Backup: history}

\begin{itemize}
\item \textbf{1982:} Koide observes
$\sum m_\ell / (\sum\sqrt{m_\ell})^2 = 2/3$ for $(e,\mu,\tau)$

\item \textbf{1966:} Miyazawa proposes hadronic supersymmetry
(meson--baryon)

\item \textbf{1985:} Catto \& G\"ursey develop algebraic
hadronic SUSY

\item \textbf{2005:} Rivero: scalar composites of 5 quarks
fill $\mathrm{SO}(32)$ adjoint; sBootstrap programme begins

\item \textbf{2024:} Published in EPJC: Diophantine uniqueness,
$\mathrm{SU}(5)$ flavour, $\mathrm{SO}(32)$ embedding

\item \textbf{2026:} This work: explicit Lagrangian,
O'Raifeartaigh derivation of the identity, mass chain,
seven parameters eliminated
\end{itemize}

\end{frame}

%-------------------------------------------------------------
\begin{frame}{Backup: why $\mathrm{SO}(32)$ and not $E_8 \times E_8$?}

The symmetric $\mathbf{15} = S^2(\mathbf{5})$ is the key test.

\bigskip
\textbf{$\mathrm{SO}(32)$:} $\;\mathbf{16} = (\mathbf{5},\mathbf{3})
\oplus (\mathbf{1},\mathbf{1})$ \quad ($5\times 3 + 1 = 16$)

$\wedge^2[(\mathbf{5},\mathbf{3})]
= \underbrace{(S^2\mathbf{5},\wedge^2\mathbf{3})}_{(\mathbf{15},\bar{\mathbf{3}})}
\oplus (\wedge^2\mathbf{5}, S^2\mathbf{3})$

The $\mathbf{15}$ appears naturally: overall antisymmetry of
the composite index forces symmetric flavour $\times$
antisymmetric colour.

\bigskip
\textbf{$E_8 \times E_8$:} Every $E_8$ decomposition under
$\mathrm{SU}(5)$ yields only exterior powers
($\mathbf{5}$, $\overline{\mathbf{10}}$, $\mathbf{24}$, \ldots).

$S^2(\mathbf{5}) = \mathbf{15}$ is \textbf{algebraically forbidden}.

\bigskip
$\Longrightarrow$ If the sBootstrap is correct,
string theory must be Type~I / heterotic $\mathrm{SO}(32)$.

\end{frame}

%-------------------------------------------------------------
\begin{frame}{Backup: a Type~I construction}

Type~IIB on $T^6/\mathbb{Z}_3$ with orientifold~$\Omega$.

\medskip
Chan-Paton embedding on 16 D9-brane pairs:
\begin{equation*}
\gamma_\theta = \mathrm{diag}(\omega\cdot\mathbb{1}_5,\;
  \omega^2\cdot\mathbb{1}_5,\; \mathbb{1}_3,\; \mathbb{1}_3)
\qquad (\omega = e^{2\pi i/3})
\end{equation*}

Twisted tadpole: $\mathrm{Tr}(\gamma_\theta)
= 5\omega + 5\omega^2 + 6 = 1$ \quad\checkmark

\bigskip
\begin{columns}[T]
\column{0.5\textwidth}
\textbf{What works:}
\begin{itemize}
\item Gauge group $\mathrm{SU}(5)_f \times \mathrm{SU}(3)_c$
\item Three chiral generations from $\mathbb{Z}_3$
\item $5\times 3 + 1 = 16$ is self-referential
\item $\mathrm{SU}(3)$ root lattice: minimal triangle
  area $\sqrt{3}/4$ \\ (candidate for $gv/m = \sqrt{3}$)
\end{itemize}

\column{0.5\textwidth}
\textbf{Open:}
\begin{itemize}
\item Chiral spectrum gives antisymmetric $\mathbf{10}$,
  not symmetric $\mathbf{15}$
\item String scale $M_s \sim 1\GeV$ requires
  $\geq 3$ large extra dimensions ($R \sim 0.2\,\mu$m)
\item Full tadpole cancellation not verified
\item Orientifold sign for Sp-type projection
\end{itemize}
\end{columns}

\end{frame}

%-------------------------------------------------------------
\begin{frame}{Backup: PMNS and neutrino mixing}

The theory provides \textbf{one} neutrino mass mechanism:

\medskip
\begin{itemize}
\item $\mathrm{SU}(2)$ sector has 4 fermions:
$\psi_{X_L}$ (goldstino), $\psi_{L_2}$ (pseudo-modulus),
$\psi_{L_1,L_3}$ ($\to \mu, \tau$)
\item $\psi_{L_2}$ gets CW mass $\sim 50\,\text{meV}$
from $W_{\rm dyn} = L_1 L_2 L_3/\Lambda_L^3$
\end{itemize}

\bigskip
\textbf{Why large PMNS angles are natural here:}

\medskip
The $s$-confining superpotential $L_1 L_2 L_3/\Lambda_L^3$
is \textbf{$S_3$-symmetric} --- it treats all three composites
democratically.

\medskip
Contrast with quarks: the Seiberg seesaw $\langle M_j\rangle = C/m_j$
creates a \emph{hierarchical} texture $\to$ small CKM angles.

\medskip
\textbf{Same Lagrangian} explains both:
\begin{center}
\begin{tabular}{lll}
\toprule
Sector & Confinement & Mixing \\
\midrule
Quarks ($\mathrm{SU}(3)$) & Seiberg seesaw (hierarchical)
  & small CKM \\
Leptons ($\mathrm{SU}(2)$) & $s$-confinement (democratic)
  & large PMNS \\
\bottomrule
\end{tabular}
\end{center}

\end{frame}

%-------------------------------------------------------------
\begin{frame}{Backup: top--neutrino duality}

\textbf{Conjecture}: the top and neutrino are ``dual'' particles ---
one gets mass in each frame of the electric-magnetic duality.

\bigskip
\begin{center}
\begin{tabular}{lll}
\toprule
Frame & Top & Neutrino \\
\midrule
Electric & Dirac mass ($y_t \sim 1$) & massless \\
Magnetic & massless & Majorana ($\sim (m/\Lambda_L)^6$) \\
\bottomrule
\end{tabular}
\end{center}

\bigskip
Both sit \textbf{outside} the main composite spectrum:
\begin{itemize}
\item Top = Chan-Paton singlet (the ``$1$'' in $5\times 3 + 1 = 16$)
\item Neutrino = pseudo-modulus of SU(2) $s$-confinement
\end{itemize}

\medskip
$D{=}11$ SUGRA: $128 = 44 + 84$, graviton $44 = 32 + 12$.

The ``$12$'' = top (electric) or neutrino (magnetic).

\medskip
The product $m_t \times m_\nu$ should be set by the duality
scale --- determined by M-theory geometry.

\end{frame}

%-------------------------------------------------------------
\begin{frame}{Backup: is this gravity?}

A string at $M_s \sim 1\GeV$ is also a theory of gravity.
Closed strings $\to$ gravitons.

\bigskip
\textbf{The puzzle is inverted:} not ``why is gravity so weak?''
but ``why is $M_{\rm Pl}$ so large compared to the hadronic scale?''

\medskip
Answer (ADD): extra dimensions.
$M_{\rm Pl}^2 \sim M_s^{n+2} R^n / g_s^2$

\medskip
\begin{center}
\begin{tabular}{ccl}
\toprule
Large dims $n$ & $R$ & Status \\
\midrule
4 & $\sim 1\;\mu$m & marginally safe \\
5 & $\sim 12$~nm & safe \\
6 & $\sim 0.6$~nm & safe \\
\bottomrule
\end{tabular}
\end{center}

\bigskip
\textbf{What changes:}
\begin{itemize}
\item Gravity becomes strong at $\sim 1\GeV$, not $10^{19}\GeV$
\item KK graviton tower starts at $1/R \sim$~keV--MeV
\item Moduli stabilisation of the orbifold $\to$ \emph{could} fix
  the pseudo-modulus at $\sqrt{3}$
\end{itemize}

\medskip
We do not claim this is gravity. But if $SO(32)$ is literal,
it comes with gravity for free.

\end{frame}

\setcounter{framenumber}{\value{finalframe}}

\end{document}
